% GENERATED FILE. Edit canonical lessons, then run npm run generate:books.
% canonical-source-hash: fnv1a64:e841b2248e751b2a
% canonical-lessons: TE-C305-yap, TE-C305-kibordu, TE-C305-lyaptap, TE-C305-byatari, TE-C305-vidiyo

\chapter{App, Keyboard, Laptop, Battery, Video}
\label{ch:te-app-keyboard-laptop-battery-video}
\hlchaptermodality{305}

\begin{chapteropening}
\textbf{By the end of this chapter:} \emph{I can say an app, a keyboard, a laptop, a battery and a video.}

Complete the last lesson of chapter 305: I can say an app, a keyboard, a laptop, a battery and a video.
\end{chapteropening}

\section[\texorpdfstring{\te{యాప్} (yāp)}{యాప్ (yāp)}]{\te{యాప్} (yāp) — an app}

\label{lesson:TE-C305-yap}



\begin{quote}
Before the new one: say the Telugu for cricket, then the Telugu for chess.
\end{quote}

\subsection*{You'll want to know: \te{యాప్}}
\textbf{\te{యాప్}} — \emph{yāp} — \textquotedblleft{}an app\textquotedblright{}.

\begin{grammarlens}[title={What you've built}]
One more word for this chapter.
\end{grammarlens}

\subsection*{Your turn}
\begin{itemize}
\raggedright
  \item \emph{Say it:} \emph{yāp}
  \item \emph{Say it:} \emph{yāp}, once more
  \item \emph{Say it:} say \emph{yāp}
  \item \emph{Recall:} say the Telugu for cricket, then the Telugu for chess, then say \emph{yāp} again
\end{itemize}

\begin{tcolorbox}[breakable,colback=teal!4,colframe=teal!35!black,title={Before you move on}]
What does \te{యాప్} mean? (\textquotedblleft{}An app\textquotedblright{}.) Say it once more.
\end{tcolorbox}
\section[\texorpdfstring{\te{కీబోర్డు} (kībōrḍu)}{కీబోర్డు (kībōrḍu)}]{\te{కీబోర్డు} (kībōrḍu) — a keyboard}

\label{lesson:TE-C305-kibordu}



\begin{quote}
Before the new one: say the Telugu for chess, then the Telugu for an app.
\end{quote}

\subsection*{You'll want to know: \te{కీబోర్డు}}
\textbf{\te{కీబోర్డు}} — \emph{kībōrḍu} — \textquotedblleft{}a keyboard\textquotedblright{}.

\begin{grammarlens}[title={What you've built}]
One more word for this chapter.
\end{grammarlens}

\subsection*{Your turn}
\begin{itemize}
\raggedright
  \item \emph{Say it:} \emph{kībōrḍu}
  \item \emph{Say it:} \emph{kībōrḍu}, once more
  \item \emph{Say it:} say \emph{kībōrḍu}
  \item \emph{Recall:} say the Telugu for chess, then the Telugu for an app, then say \emph{kībōrḍu} again
\end{itemize}

\begin{tcolorbox}[breakable,colback=teal!4,colframe=teal!35!black,title={Before you move on}]
What does \te{కీబోర్డు} mean? (\textquotedblleft{}A keyboard\textquotedblright{}.) Say it once more.
\end{tcolorbox}
\section[\texorpdfstring{\te{ల్యాప్టాప్} (lyāpṭāp)}{ల్యాప్టాప్ (lyāpṭāp)}]{\te{ల్యాప్టాప్} (lyāpṭāp) — a laptop}

\label{lesson:TE-C305-lyaptap}



\begin{quote}
Before the new one: say the Telugu for an app, then the Telugu for a keyboard.
\end{quote}

\subsection*{You'll want to know: \te{ల్యాప్టాప్}}
\textbf{\te{ల్యాప్టాప్}} — \emph{lyāpṭāp} — \textquotedblleft{}a laptop\textquotedblright{}.

\begin{grammarlens}[title={What you've built}]
One more word for this chapter.
\end{grammarlens}

\subsection*{Your turn}
\begin{itemize}
\raggedright
  \item \emph{Say it:} \emph{lyāpṭāp}
  \item \emph{Say it:} \emph{lyāpṭāp}, once more
  \item \emph{Say it:} say \emph{lyāpṭāp}
  \item \emph{Recall:} say the Telugu for an app, then the Telugu for a keyboard, then say \emph{lyāpṭāp} again
\end{itemize}

\begin{tcolorbox}[breakable,colback=teal!4,colframe=teal!35!black,title={Before you move on}]
What does \te{ల్యాప్టాప్} mean? (\textquotedblleft{}A laptop\textquotedblright{}.) Say it once more.
\end{tcolorbox}
\section[\texorpdfstring{\te{బ్యాటరీ} (byāṭarī)}{బ్యాటరీ (byāṭarī)}]{\te{బ్యాటరీ} (byāṭarī) — a battery}

\label{lesson:TE-C305-byatari}



\begin{quote}
Before the new one: say the Telugu for a keyboard, then the Telugu for a laptop.
\end{quote}

\subsection*{You'll want to know: \te{బ్యాటరీ}}
\textbf{\te{బ్యాటరీ}} — \emph{byāṭarī} — \textquotedblleft{}a battery\textquotedblright{}.

\begin{grammarlens}[title={What you've built}]
One more word for this chapter.
\end{grammarlens}

\subsection*{Your turn}
\begin{itemize}
\raggedright
  \item \emph{Say it:} \emph{byāṭarī}
  \item \emph{Say it:} \emph{byāṭarī}, once more
  \item \emph{Say it:} say \emph{byāṭarī}
  \item \emph{Recall:} say the Telugu for a keyboard, then the Telugu for a laptop, then say \emph{byāṭarī} again
\end{itemize}

\begin{tcolorbox}[breakable,colback=teal!4,colframe=teal!35!black,title={Before you move on}]
What does \te{బ్యాటరీ} mean? (\textquotedblleft{}A battery\textquotedblright{}.) Say it once more.
\end{tcolorbox}
\section[\texorpdfstring{\te{వీడియో} (vīḍiyō)}{వీడియో (vīḍiyō)}]{\te{వీడియో} (vīḍiyō) — a video}

\label{lesson:TE-C305-vidiyo}



\begin{quote}
Before the new one: say the Telugu for a laptop, then the Telugu for a battery.
\end{quote}

\subsection*{You'll want to know: \te{వీడియో}}
\textbf{\te{వీడియో}} — \emph{vīḍiyō} — \textquotedblleft{}a video\textquotedblright{}.

\begin{grammarlens}[title={What you've built}]
One more word for this chapter.
\end{grammarlens}

\subsection*{Your turn}
\begin{itemize}
\raggedright
  \item \emph{Say it:} \emph{vīḍiyō}
  \item \emph{Say it:} \emph{vīḍiyō}, once more
  \item \emph{Say it:} say \emph{vīḍiyō}
  \item \emph{Recall:} say the Telugu for a laptop, then the Telugu for a battery, then say \emph{vīḍiyō} again
\end{itemize}

\begin{tcolorbox}[breakable,colback=teal!4,colframe=teal!35!black,title={Before you move on}]
What does \te{వీడియో} mean? (\textquotedblleft{}A video\textquotedblright{}.) Say it once more.
\end{tcolorbox}
